\section{Persistencia, estructuras de fusión y realizaciones espectrales}
\label{vi:sec:topologia}

El atlas y los operadores de masa admiten distintas representaciones. La
clasificación de esas representaciones conserva el dato que caracteriza
cada una: una regla de multiplicación, una conjugación real, una fase de
intercambio o una medida espectral. Esta sección sitúa los sectores
topológicos y las realizaciones neutras dentro de la misma arquitectura.
Las coordenadas de masa se obtienen una vez que la representación y su
dinámica están determinadas.

La diferencia entre estas etapas es especialmente visible en el sector
areal--volumétrico. Una matriz de incidencia puede fijar un anillo de fusión
y su dimensión positiva; la energía de una excitación depende también del
operador que actúa en la representación. La primera determinación conserva
un contenido exacto cuando se estudian varias realizaciones dinámicas
del mismo anillo. La distinción mantiene la información estructural que
precede a la magnitud.

\subsection{El anillo basado de la incidencia areal--volumétrica}
\label{vi:subsec:fusion-av}

La incidencia primitiva de las hojas, obtenida del calendario del TPK, es
\[
 F_{\mathrm{av}}=\begin{pmatrix}0&1\\1&1\end{pmatrix}.
\]
En la base $(1,\tau)$, la primera columna expresa
$\tau\cdot1=\tau$; la otra expresa $\tau\cdot\tau=1+\tau$.
La incidencia induce así el anillo basado
\[
 \mathcal F_{\mathrm{av}}=\mathbb Z[\tau]/(\tau^2-\tau-1),
 \qquad \mathcal F_{\mathrm{av}}^+
 =\{a+b\tau:a,b\in\mathbb Z_{\geq0}\}.
\]
El cono positivo registra multiplicidades. La operación es interna a los
dos tipos y no utiliza una masa como coeficiente.

\begin{proposition}[Producto y dimensión positiva]
\label{vi:prop:fusion-av}
Se cumple
\[
 (a+b\tau)(c+d\tau)
 =(ac+bd)+(ad+bc+bd)\tau.
\]
Existe un único homomorfismo de anillos $d:\mathcal F_{\mathrm{av}}
\to\mathbb R$ con $d(1)=1$ y $d(\tau)>0$. Su valor en $\tau$
es la coordenada áurea reconocida en la realización de $F_{\mathrm{av}}$.
\end{proposition}
\begin{proof}
Expandir el producto y sustituir $\tau^2=1+\tau$ da los coeficientes.
Un homomorfismo unitario queda determinado por $x=d(\tau)$, con
$x^2=x+1$. Sus raíces son $(1+\sqrt5)/2$ y $(1-\sqrt5)/2$;
sólo la primera es positiva. La evaluación en esa raíz define un
homomorfismo del cociente. La ecuación identifica la dimensión de esta
representación; la coordenada regional HMT ha sido construida antes.
\end{proof}

Con $F_0=0$, $F_1=1$ y $F_{n+1}=F_n+F_{n-1}$, las potencias satisfacen
\[
 \tau^n=F_{n-1}+F_n\tau\qquad(n\geq1).
\]
El caso inicial es inmediato; multiplicar por $\tau$ aplica de nuevo
la incidencia y produce el paso inductivo. Los coeficientes conservan
así su origen como recuentos de composición de las hojas.

Una representación categórica añade espacios de morfismos y un
asociador a estas multiplicidades. Una representación trenzada añade
operadores de intercambio compatibles. La realización material añade
el Hamiltoniano, sus dominios y la preparación que selecciona un
sector espectral. Si posee brecha $\Delta E>0$ y velocidad efectiva
$c_{\mathrm{eff}}$, su coordenada de masa de brecha es
\[
 m_{\mathrm{brecha}}=\Delta E\,c_{\mathrm{eff}}^{-2}.
\]
La dimensión positiva del anillo y la brecha son salidas de operadores
diferentes. La regla de fusión conserva su contenido sin fijar, por
ella sola, una masa universal para todas sus realizaciones.

\subsection{Dimensión y tipo de representación}

En una representación apuntada, todos los objetos simples son
invertibles. Si $X$ es simple invertible y $d$ es una dimensión
positiva compatible con duales, entonces
\[
 1=d(X\otimes X^*)=d(X)d(X^*)=d(X)^2,
\]
y $d(X)=1$. La dimensión proporciona, por tanto, un control de los
cambios de representación.

\begin{proposition}[Obstrucción de dimensión en el sector apuntado]
\label{vi:prop:apuntada}
Un simple de dimensión $\varphi$ o $\sqrt2$ no puede identificarse,
preservando la dimensión, con un simple de una representación apuntada.
En particular, los $81$ simples invertibles del centro no torcido
basado en $\mathbb Z/9\mathbb Z$ no proporcionan esa identificación.
\end{proposition}
\begin{proof}
Los simples invertibles tienen dimensión $1$, distinta de ambas
dimensiones. En el centro no torcido abeliano indicado, un simple se
etiqueta por un elemento del grupo y un carácter de su dual. Ambos
conjuntos tienen $9$ elementos. Las $81$ etiquetas son invertibles
por la multiplicación del grupo y del carácter.
\end{proof}

Este censo $81$ tiene un dominio distinto del atlas de $81$ celdas.
Una extensión o un funtor hacia otro codominio declara las acciones
y los mapas que transporta. La incidencia areal--volumétrica conserva
su anillo exacto y la realización topológica conserva las condiciones
de su propio tipo.

\subsection{Autoconjugación neutra y generadores reales de CAR}
\label{vi:subsec:majorana-categorias}

Una estructura real en un espacio complejo es un operador antiunitario
$\mathcal C$ con $\mathcal C^2=I$. Su parte fija es un espacio real.
La conmutación del operador de masa con $\mathcal C$ preserva esa
parte real en su dominio. La restricción de una evolución a la parte
fija se expresa mediante la conmutación de su grupo con $\mathcal C$;
en una evolución de Schr\"odinger hay que distinguir esta condición
de la conmutación del Hamiltoniano. Ambas condiciones pertenecen a la
representación material de la sección~\ref{vi:subsec:antisecciones}.

\begin{proposition}[Conjugación antiunitaria y evolución temporal]
\label{vi:prop:conjugacion-evolucion}
Sean $H$ autoadjunto, $\hbar>0$ y $\mathcal C$ una conjugación
antiunitaria que conserva $\mathcal D(H)$. Para
$U(t)=\exp(-itH/\hbar)$, si
$\mathcal C H\mathcal C^{-1}=\epsilon H$, con
$\epsilon\in\{1,-1\}$, se cumple
\begin{equation}
 \mathcal C U(t)\mathcal C^{-1}=U(-\epsilon t).
 \label{vi:eq:conjugacion-grupo-temporal}
\end{equation}
El grupo $U(t)$ preserva la parte fija de $\mathcal C$ para todo $t$
si y sólo si conmuta con $\mathcal C$, equivalentemente si
$\mathcal C H\mathcal C^{-1}=-H$.
Si, en cambio, $\mathcal C H=H\mathcal C$, un vector fijo por
$\mathcal C$ permanece fijo por ella durante toda la evolución
exactamente cuando pertenece a $\ker H$.
\end{proposition}
\begin{proof}
La conjugación antiunitaria envía $i$ a $-i$. El cálculo funcional
espectral proporciona
\[
 \mathcal C\exp(-itH/\hbar)\mathcal C^{-1}
 =\exp\bigl(it\mathcal C H\mathcal C^{-1}/\hbar\bigr),
\]
que prueba \eqref{vi:eq:conjugacion-grupo-temporal}. Si $U(t)$ conmuta
con $\mathcal C$, transporta sus vectores fijos a vectores fijos.
Recíprocamente, todo vector complejo se escribe de manera única como
$x+iy$ con $x,y$ fijos por $\mathcal C$. La preservación de esa parte
real implica, por linealidad de $U(t)$ y antilinealidad de
$\mathcal C$, que ambos conmutan. La unicidad del generador
autoadjunto del grupo unitario identifica entonces
$\mathcal C H\mathcal C^{-1}$ con $-H$.

En el caso conmutativo $\epsilon=1$, para $\mathcal C\psi=\psi$
se tiene $\mathcal C U(t)\psi=U(-t)\psi$. Que $U(t)\psi$ sea fijo
para todo $t$ equivale a $U(2t)\psi=\psi$ para todo $t$. Por el
cálculo espectral, esto equivale a que la medida espectral de $\psi$
esté soportada en $\{0\}$, es decir, a $\psi\in\ker H$.
\end{proof}

Así, la simetría antiunitaria que conmuta con $H$ invierte el sentido
temporal. La realización mediante estados reales conservados utiliza
otra condición sobre el generador. Esta distinción preserva el alcance
de la autoconjugación algebraica y permite especificar por separado la
dinámica de un campo de Majorana.

\begin{proposition}[Autoconjugación y carga]
\label{vi:prop:carga-real}
Sean $Q$ un operador de carga autoadjunto y $\mathcal C$ una
conjugación que transporta su dominio y satisface
$\mathcal C Q\mathcal C^{-1}=-Q$. Si $\psi\ne0$,
$Q\psi=q\psi$ y $\mathcal C\psi=\psi$, entonces $q=0$.
\end{proposition}
\begin{proof}
La autoadjunción implica $q\in\mathbb R$. Aplicar $\mathcal C$
a la ecuación propia da $-Q\psi=q\psi$. Con la ecuación inicial
resulta $2q\psi=0$, luego $q=0$.
\end{proof}

La neutralidad es necesaria en este problema de autoconjugación;
el núcleo de $Q$ no especifica, por sí solo, un operador antiunitario
ni su compatibilidad con el grupo de evolución. El criterio se aplica a los
bloques neutrínicos después de construir sus operadores y su mezcla.
La palabra dodecafásica fija, el espinor real y la representación de
carga conservan así sus respectivos dominios.

La completación fermiónica proporciona otra realización. En el
espacio exterior de los modos declarados, $a_j^*$ es producto
exterior por el vector básico y $a_j$ su adjunto. La antisimetría
exterior da
\[
 \{a_j,a_k\}=0,\qquad \{a_j,a_k^*\}=\delta_{jk}I.
\]

\begin{proposition}[Generadores autoadjuntos de CAR]
\label{vi:prop:car-real}
Los operadores
\[
 \gamma_{2j-1}=a_j+a_j^*,\qquad
 \gamma_{2j}=-i(a_j-a_j^*)
\]
son autoadjuntos y cumplen $\{\gamma_r,\gamma_s\}=2\delta_{rs}I$.
\end{proposition}
\begin{proof}
La adjunción deja fijo cada operador. Para el mismo modo,
$a_j^2=(a_j^*)^2=0$ y $a_ja_j^*+a_j^*a_j=I$, luego cada
cuadrado vale $I$. El anticomutador entre las dos combinaciones
del mismo modo se anula por expansión. Para modos distintos,
todos los anticomutadores elementales se anulan.
\end{proof}

Estos son generadores reales de una estructura de Clifford. Un modo
cero topológico requiere además que la dinámica realice una excitación
localizada de energía cero, separada del continuo por una brecha y
estable bajo las perturbaciones consideradas. La propiedad algebraica
y la espectral quedan vinculadas por la realización dinámica.

\subsection{Fases de intercambio y devanado de orden seis}

Para \(n\ge2\), en el lector de orden seis la suma de exponentes define
$w_6:B_n\to\mathbb Z/6\mathbb Z$. Las relaciones de Artin
conservan esa suma: la conmutación de generadores distantes mantiene
dos términos y la relación de tres generadores conserva tres.

\begin{proposition}[Caracteres del lector de devanado]
\label{vi:prop:devanado-seis}
Los caracteres unitarios del cociente son
\[
 \chi_k(b)=\exp\!\left(\frac{2\pi i k}{6}w_6(b)\right),
 \qquad k=0,1,\ldots,5.
\]
Una representación escalar de esta forma desciende al grupo de
permutaciones exactamente para $k=0$ o $k=3$.
\end{proposition}
\begin{proof}
Un carácter cíclico queda determinado por la imagen $q$ de $1$,
con $q^6=1$. Para descender a permutaciones debe satisfacer la
relación adicional $\sigma_j^2=1$, que exige $q^2=1$.
Entre las raíces sextas sólo $q=1,-1$ cumplen esa condición.
Ambas satisfacen todas las relaciones del grupo de permutaciones.
\end{proof}

Los otros cuatro caracteres conservan información de devanado.
Una especie material necesita que el lector se realice sobre una
multisección y que la dinámica conserve su estructura. La inversión
de ruta cambia $w_6$ por $-w_6$ y transporta cada carácter a su
conjugado complejo.

\subsection{Estabilidad y visibilidad de una sección}

La positividad demostrada en la
sección~\ref{vi:sec:operador-masa} tiene una consecuencia espectral.
Para un operador autoadjunto positivo $M$, el cálculo funcional da
$\operatorname{Spec}(M^2)\subset[0,\infty)$. En una realización que
identifica $M^2$ con el observable de masa al cuadrado sobre un
espacio físico de producto positivo, una sección propia persistente
no puede tener autovalor negativo de ese observable. Una dirección
inestable de la linealización de un fondo pertenece a otro problema
espectral, que se estudia con su dominio y su dinámica.

La visibilidad depende del acoplamiento. Si $P_X$ es el soporte de
una sección y $P_a$ el proyector de un canal, $P_aP_X=0$ expresa
ortogonalidad de soportes. El producto no nulo expresa incidencia
posible; su intensidad la determina el operador de acoplamiento.
Dos proyectores pueden conmutar y tener producto no nulo, por lo
que el conmutador no sustituye esa incidencia.

En un bloque neutro, una sección estéril respecto del canal débil
pertenece al núcleo de su acoplamiento. Si ese núcleo es invariante
para el operador de masa, la restricción posee una medida espectral
propia. Profundidad del atlas, neutralidad y ausencia de acoplamiento
débil son propiedades distintas que se verifican sobre la misma
sección. Las posibilidades representativas se conservan sin
convertir una etiqueta de colección en una masa externa.

La organización mantiene así el orden del argumento: la incidencia
determina multiplicidades; los levantamientos determinan
conjugaciones y fases; la representación y el operador material
determinan la salida espectral. El contraste se aplica después a
esa salida, con sus unidades y su acoplamiento declarados.
